
Comme les deux membres de la relation à démontrer sont additifs, on peut supposer que \(x\) est la classe d'un \(\underline{O}_Y\)-Module localement libre de type fini \(\underline{M}\). L'élément \(i^*(i_*(x))\) est alors la classe du complexe \(\underline{\underline{L}} i^*(\underline{\underline{R}}i_*(\underline{M}))\), obtenu en prenant une résolution projective du complexe \(\underline{\underline{R}}i_*(\underline{M})=i_*(\underline{M})\), et en la tensorisant par \(\underline{O}_Y\). Les faisceaux d'homologie de ce complexe sont donc les \(\underline{\Tor}_i^{\underline{O}_X}(\underline{M},\underline{O}_Y)\). D'après 2.4, on a l'isomorphisme:
\[
\underline{\Tor}_i^{\underline{O}_X}(\underline{M},\underline{O}_Y)
\xrightarrow{\sim}
\underline{M}\otimes_{\underline{O}_Y}\underline{\Tor}_i^{\underline{O}_X}(\underline{O}_Y,\underline{O}_Y).
\]
D'après 2.5, on a l'isomorphisme:
\[
\underline{\Tor}_i^{\underline{O}_X}(\underline{O}_Y,\underline{O}_Y)
\xrightarrow{\sim}
\bigwedge^i\underline{N}.
\]

Les \(\underline{\Tor}_i^{\underline{O}_X}(\underline{M},\underline{O}_Y)\) sont donc des \(\underline{O}_Y\)-Modules localement libres de type fini, et la classe dans \(K^{\bullet}(Y)\) du complexe \(\underline{\underline{L}} i^*(\underline{\underline{R}}i_*(\underline{M}))\) est la somme alternée des classes de ses faisceaux d'homologie, soit:
\[
i^*(i_*(x))=\sum_{i\ge0}(-1)^i
\operatorname{cl}\bigl(\underline{M}\otimes_{\underline{O}_Y}\underline{\Tor}_i^{\underline{O}_X}(\underline{O}_Y,\underline{O}_Y)\bigr)
=\sum_{i\ge0}(-1)^i x.\lambda^i(N).
\]

\textbf{Corollaire 2.8.} \uline{Sous les hypothèses de} 2.7, \uline{soient} \(x,y\in K^{\bullet}(Y)\); \uline{on a la relation}:
\begin{equation}\notag
i_*(xy\lambda_{-1}(N))=i_*(x).i_*(y).
\end{equation}

D'après 2.7, \(y\lambda_{-1}(N)=i^*(i_*(y))\).

Par application de la formule de projection, on trouve donc:
\[
i_*(xy\lambda_{-1}(N))
=i_*(x.i^*(i_*(y)))
=i_*(x).i_*(y).
\]

\section*{3. Calcul du \texorpdfstring{$K^{\bullet}$}{K} d'un schéma éclaté}

\textbf{3.1.} Dans ce paragraphe, on désigne par \(i:Y\to X\) une immersion fermée régulière. On note \(X'\) le schéma obtenu en faisant éclater \(Y\), et \(Y'=Y\times_X X'\).
